%----------------------------------------------------------------------------------------
% ONTBINDEN IN FACTOREN. Alle opgaven zijn eigen werk.
%
% Een hoofdstuk kernvaardigheden (ontwerpregels, deel Kernvaardigheden): beknopte theorie, daarna drill.
% Labels van dit hoofdstuk beginnen met h3-.
%----------------------------------------------------------------------------------------

\mychapter{3}{Ontbinden in factoren}

\lkr{Herhaling van de 1ste graad (afzonderen, merkwaardige producten) en van Tweedegraadsfuncties en Veeltermen (som-product, discriminant, schema van Horner). Bundel 5 (groeperen) en 6 (Horner) veronderstellen Veeltermen; sla ze over als dat handboek nog niet gezien is. Nodig voor breuken vereenvoudigen (hoofdstuk 4) en vergelijkingen oplossen met de product-nulregel (hoofdstuk 5).}

\section{Wat je moet kennen}

\begin{theorie}
\begin{minipage}[t]{.47\linewidth}
\vspace{0pt}
$\formule{A^2 - B^2 = (A - B)(A + B)}$

\vspace{1mm}
$\formule{A^2 \pm 2AB + B^2 = (A \pm B)^2}$

\vspace{1mm}
$\formule{x^2 + bx + c = (x + p)(x + q)}$ \\ \hspace*{6mm} met $p + q = b$ en $p \cdot q = c$
\end{minipage}\hfill
\begin{minipage}[t]{.5\linewidth}
\vspace{0pt}
Heeft $ax^2 + bx + c$ met $\formule{D = b^2 - 4ac \geqslant 0}$ de oplossingen $x_1$ en $x_2$, dan is

\vspace{1mm}
$\formule{ax^2 + bx + c = a(x - x_1)(x - x_2)}$

\vspace{1mm}
Is $D < 0$, dan kun je niet ontbinden in factoren van graad $1$.
\end{minipage}
\end{theorie}

\begin{werkwijze}
\textbf{1.} Breng een gemeenschappelijke factor buiten haakjes. \ \textbf{2.} Twee termen: een verschil van kwadraten? Drie termen: een volkomen kwadraat, de som-productmethode of de discriminant. Vier termen: groeperen. Hogere graad: zoek een nulwaarde $a$ bij de delers van de constante term en deel door $x - a$ met het schema van Horner. \ \textbf{3.} Kijk bij elke factor of je verder kunt ontbinden.
\end{werkwijze}

\textbf{Klassieke fouten.} Niet verder ontbinden: $2x^2 - 8 = 2(x^2 - 4)$ is niet af, want $x^2 - 4 = (x - 2)(x + 2)$. Een som van kwadraten ontbinden: $x^2 + 4$ is \emph{niet} $(x + 2)^2$, en ontbindt niet. Bij $a \neq 1$ de factor $a$ vergeten: $2x^2 - x - 3 = 2\bigl(x - \tfrac32\bigr)(x + 1) = (2x - 3)(x + 1)$.

\voorbeeld{}{Ontbind volledig: \quad a. $3x^3 - 12x$ \qquad b. $x^3 - 7x + 6$\label{vb:h3-ontbind}}{%
\naast{\begin{stappenvrij}{6.6cm}{4.2cm}
3x(x^2 - 4) = 3x(x - 2)(x + 2) & a. Eerst $3x$ afzonderen, dan $A^2 - B^2$. \\
1 - 7 + 6 = 0 & b. $1$ is een nulwaarde: deel door $x - 1$. \\
x^2 + x - 6 = (x - 2)(x + 3) & Het quotiënt met som-product. \\
x^3 - 7x + 6 = (x - 1)(x - 2)(x + 3)
\end{stappenvrij}}{\horner{3}{1}{1 & 0 & -7 & 6}{ & 1 & 1 & -6}{1 & 1 & -6}{0}}}

%========================================================================================
\opdrachtenkop

\begin{oefeningen}
\needspace{10\baselineskip}
\oef{h3-o1}

\needspace{10\baselineskip}
\oef{h3-o2}

\needspace{10\baselineskip}
\oef{h3-o3}

\needspace{10\baselineskip}
\oef{h3-o4}

\needspace{10\baselineskip}
\oef{h3-o5}

\needspace{10\baselineskip}
\oef{h3-o6}

\needspace{10\baselineskip}
\oef{h3-o7}

\needspace{10\baselineskip}
\oef{h3-o8}

\needspace{10\baselineskip}
\oef{h3-o9}

\needspace{10\baselineskip}
\oef{h3-o10}
\end{oefeningen}

%----------------------------------------------------------------------------------------
\needspace{14\baselineskip}
\opdrachtsectie{Test jezelf}

\begin{oefeningen}
\oef{h3-t1}

\oef{h3-t2}

\oef{h3-t3}

\oef{h3-t4}
\end{oefeningen}

\herhaalschema{%
  \fzv{oef:h3-t1}{oef:h3-o2}%
  \fzv{oef:h3-t2}{oef:h3-o4}%
  \fv{oef:h3-t3}{vb:h3-ontbind}{oef:h3-o6}%
  \fv{oef:h3-t4}{vb:h3-ontbind}{oef:h3-o7}}

\ifzelfstudie
\section*{Maak je eigen samenvatting}
Sluit dit boek en schrijf in je schrift: de stappen om te ontbinden in factoren, de merkwaardige producten in de richting van ontbinden, de som-productmethode, de ontbinding met de discriminant, en wanneer een veelterm van graad $2$ niet te ontbinden is. Vergelijk daarna met het kader.
\fi
